\section{Discussion}

\subsection{Architectural Synthesis and Theoretical Implications}
The empirical results compiled across 8,654 authentic software specifications substantiate the central hypothesis of this investigation: translating natural language requirements into synthetic biological genomic sequences through a dynamic syntactic-semantic dinucleotide codon grammar establishes an invariant, highly discriminative representation that substantially outperforms pure textual and deep language modeling baselines. From an architectural perspective, natural language requirements engineering has long been caught in a computational tension between two extremes. On one end, classical sparse lexical representations (such as bag-of-words and term frequency matrices) offer rapid, deterministic extraction but suffer from extreme feature sparsity, lack of contextual syntactic depth, and acute vulnerability to vocabulary mismatch [1,2,4]. On the other end, massive pre-trained transformer encoders capture rich semantic associations but impose prohibitive computational latencies, require heavy GPU infrastructure, and remain susceptible to domain hallucination and subword truncation when processing terse, jargon-laden software clauses [3,5].

DNA+ bridges this methodological chasm by exploiting the formal mathematical isomorphisms between macromolecular genomic sequences and human symbolic language [6,8]. By mapping functional imperative anchors and auxiliary constraint modalities into coupled nucleotide dimers ($AA, AT, \dots, GG$), the framework constructs a continuous structural backbone for each requirement statement. The subsequent extraction of alignment-free oligonucleotide statistics—including sixteen dinucleotide relative abundance ratios, GC-content, AT/GC skew, and Shannon informational entropy—yields a compact 23-dimensional physicochemical profile that functions as an invariant structural signature [7,9,10]. As demonstrated by our ablation experiments, this 23-dimensional sequence manifold in isolation achieves 93.81\% accuracy, dramatically outperforming 434 dimensions of classical NLP features (59.13\%). When fused with 384-dimensional Sentence-BERT embeddings [11] and calibrated lexical n-grams, the resulting 457-dimensional multimodal manifold captures both high-order abstract semantics and low-level syntactic regularity, unlocking peak generalization across heterogeneous industrial benchmarks.

\subsection{Robustness Under Extreme Class Distribution Imbalance}
A persistent challenge in empirical software engineering is the severe class distribution skew inherent to requirements specifications [17,20]. In authentic enterprise documentation, functional requirements naturally dominate system narratives, whereas critical non-functional attributes (such as fault tolerance, scalability, and maintainability) are sparsely distributed. On the extended PROMISE-EXP benchmark, this distributional skew reaches an extreme imbalance ratio of 37.0:1, while the industrial FNFC corpus exhibits a 5.85:1 operational disparity. Conventional machine learning frameworks confronted with such skews routinely experience severe performance degradation or collapse into majority-class triviality.

To address class skew, prevailing methodologies in requirements engineering rely heavily on synthetic data oversampling (e.g., SMOTE) or generative data augmentation via back-translation and large language models [16,27,28]. However, as established in foundational empirical software engineering literature, synthetic oversampling techniques risk introducing artificial data artifacts and distorting real-world decision boundaries [18,19]. Rather than altering raw empirical distributions, DNA+ addresses class imbalance natively through two complementary mechanisms. First, the dinucleotide codon grammar provides orthogonal topological separation: because non-functional attributes possess unique constraint structures (reflected in high GC-content and distinct transition odds), rare classes occupy distinct geometric clusters within the sequence manifold. Second, cost-sensitive class re-weighting penalizes minority misclassifications directly during ensemble training. This dual strategy is empirically substantiated by our Precision-Recall calibration analysis, where Gradient Boosting attained a macro-averaged PR-AUC of 0.9361 across twelve classes, proving that minority categories are preserved with high precision without generating synthetic text artifacts.

\subsection{Comparative Superiority Over 2025--2026 State-of-the-Art Paradigms}
When evaluated against the contemporary state-of-the-art literature published in 2025 and 2026, DNA+ establishes a decisive performance advantage ranging from +5.17\% to +9.67\% in classification accuracy and +5.50\% to +10.30\% in macro-F1 score. A rigorous comparative examination reveals the specific methodological factors driving this superiority.

Rahman et al. [26] developed RequiBERT, a specialized transformer model fine-tuned on software requirement specifications, achieving 89.60\% accuracy and 88.40\% macro-F1. While domain-specific pre-training enhances contextual comprehension, pure transformer encoders remain prone to attention dilution on brief clauses and incur severe computational overhead during inference. DNA+ surpasses RequiBERT by +5.17\% (94.77\% accuracy) because its biological sequence branch captures deterministic morphological cues that attention mechanisms smooth over. Naseer et al. [27] tackled requirement sparsity through multilingual back-translation coupled with BERT (89.20\% accuracy, 87.90\% macro-F1). However, back-translation frequently introduces semantic drift and subtle grammatical alterations into technical constraints. DNA+ outperforms their framework by +5.57\% while remaining entirely deterministic and hallucination-free.

Similarly, AlDhafer et al. [28] employed rule-based lexical substitution combined with BiLSTM-RoBERTa (87.50\% accuracy), yet experienced persistent error propagation on rare architectural classes. DNA+ establishes a +7.27\% advantage, confirming that bio-inspired oligonucleotide frequencies provide invariant topological anchors that isolate rare classes naturally. In ensemble modeling, Limaylla-Lunarejo et al. [29] investigated meta-stacking of traditional classifiers over TF-IDF and GloVe vectors, reaching 86.40\% accuracy. DNA+ surpasses their meta-ensemble by +8.37\%, proving that superior feature representation over a bio-inspired multimodal manifold yields far greater performance dividends than ensembling weak learners over impoverished lexical inputs. Finally, the pre-trained transformer benchmark of Ayed et al. [30] peaked at 85.10\% accuracy across BERT-base, DistilBERT, and RoBERTa. DNA+ demonstrates a +9.67\% lead over this transformer baseline, confirming that integrating structural dinucleotide transition matrices provides crucial inductive bias that deep language models lack.

\subsection{Actionable Software Engineering Implications and Industrial Practice}
The practical software engineering implications of the DNA+ framework extend beyond benchmark accuracy metrics to core lifecycle activities:
First, automated quality attribute triage during early elicitation: In rapid agile and continuous delivery environments, incoming user stories and technical backlogs can be processed through the DNA+ pipeline in real time. Because feature extraction relies on alignment-free oligonucleotide counts and lightweight Siamese embeddings, inference latency requires only fractions of a millisecond per requirement, enabling seamless integration into continuous integration pipelines, issue tracking systems, and automated requirements management suites.
Second, verified human-in-the-loop auditability: A persistent barrier to adopting deep learning in safety-critical domains (such as aerospace, healthcare, and automotive systems) is the opaque nature of neural predictions. Through its dual-tier explainability framework, DNA+ couples global TreeSHAP attribution with local LIME token justifications [21,22]. Requirements engineers receive transparent rationales highlighting specific linguistic anchors and sequence properties driving each classification, fostering trust and streamlining contractual compliance verification.
Third, cross-domain architectural consistency: The observed generalization across three distinct corpora—advancing from academic benchmarks (PROMISE) to large-scale commercial software repositories (FNFC, $N=7,060$)—confirms that the dynamic dinucleotide codon grammar is resilient against domain-specific vocabulary shifts, establishing a reliable foundation for cross-project transfer learning.

\subsection{Threats to Validity}
In accordance with established empirical guidelines for software engineering research [23,24,25], the potential threats to the validity of this study are categorized and addressed across four dimensions:

Threats to Internal Validity: Internal validity concerns the extent to which experimental controls prevent confounding variables from biasing observed outcomes. A primary internal threat is data leakage between training and validation partitions. To prevent information leakage, all preprocessing pipelines—including TF-IDF vocabulary fitting, SBERT normalization, and cost-sensitive weight calculations—were executed strictly within the training folds of each iteration. Furthermore, model evaluations were conducted using repeated stratified ten-fold cross-validation across thirty randomized splits (300 evaluations per algorithm), minimizing partition bias and guaranteeing that observed stability profiles (Figure 10) reflect genuine algorithmic generalization rather than favorable data splits [12,13,14,15].

Threats to External Validity: External validity relates to the generalizability of the findings across unseen software engineering domains. While requirements corpora vary in stylistic phrasing, organizational conventions, and technical domains, our evaluation spanned three distinct datasets comprising 8,654 authentic specifications ranging from canonical academic sets to enterprise-scale industrial repositories. However, all evaluated requirements were authored in English. Replicating the framework across multilingual software engineering environments, non-English specifications, or mixed-language technical documentation represents an important external boundary condition.

Threats to Construct Validity: Construct validity addresses whether the selected evaluation metrics accurately reflect real-world performance. Evaluating imbalanced software datasets solely through classification accuracy introduces substantial construct bias. To mitigate this threat, our evaluation battery paired discrete accuracy with macro-averaged F1 scores, Receiver Operating Characteristic curves, and Precision-Recall Area Under the Curve metrics across all multi-class and binary configurations, ensuring that minority-class performance was rigorously evaluated.

Threats to Conclusion Validity: Conclusion validity relates to the mathematical correctness of statistical inferences drawn from the experimental data. To guarantee conclusion validity, normality was verified via the Shapiro-Wilk test, paired parametric t-tests were corroborated with non-parametric Wilcoxon signed-rank and Mann-Whitney U tests, and inter-corpus differences were evaluated via the Friedman omnibus test. Standardized effect sizes (Cohen's $d = 4.4171$, Hedges' $g = 4.1351$, and Cliff's $\delta = 1.0000$) definitively confirmed that observed improvements represent massive, practically meaningful advances rather than minor stochastic fluctuations.
